\ssscp{{\#}muku `banana'}{02-03-01-07}
In \srf{02-01}, we noted the \it{westward} dispersal of important
food crops such as bananas,
sugarcane, taro, breadfruit and sago from the New Guinea region
prior to the arrival of the Austronesians.
The distribution of words associated with some of these items
has important implications for our understanding of the spread
of languages in the region.
\citet{DenhamDonohue2009} have looked at banana DNA
and words for ‘banana’,
and observed some interesting patterns.
Many Austronesian languages in the western
and central regions, and out into Oceania show expected reflexes
of \tsc{PAn} *punti ‘banana’.
A small sample is shown in example \qf{02-41},
(data from \citealt{BlustTrussel2020}, Grimes fieldnotes, Edwards fieldnotes).
Shading in example \qf{02-41} shows languages in our target region.

%{\st{4pt}\mbox{}\\\noindent
{\begin{xltabular}{\textwidth}{>{\hspace{-\tabcolsep}}l<{\hspace{-\tabcolsep}}>{\hspace{-\tabcolsep}\enspace}llll}
\lsptoprule
\tbx{02-41}			&	\tsc{PAn}	&	*punti	&	‘banana’	&	Region		\\ \midrule
%\endfirsthead
	%&	\tsc{PAn}	&	*punti	&	‘banana’	&	region		\\ \midrule
%\endhead
	&	Palawan	&	{\hr}punti	&	‘banana’	&	Philippines		\\
	&	Sasak	&	{\hr}puntiʔ	&	‘banana’	&	Lombok		\\
%\rcl	&	Manggarai	&	{\hr}punti	&	‘k.o. banana’	&	\tsc{Bima-Lembata}, Flores\\
%\rcl	&		&	&		&	 (Pacar dialect)		\\
\rcl	&	Tii	&	{\hr}hundi	&	‘banana’	&	\tsc{Timor-Babar}, Rote		\\
\rcl	&	Termanu	&	{\hr}huni	&	‘banana’	&	\tsc{Timor-Babar}, Rote		\\
\rcl	&	Kotos Amarasi	&	  {\hr}uki	&	‘banana’	&	\tsc{Timor-Babar}, (regular *nt > \it{k})		\\
\rcl	&	Tetun	&	{\hr}hudi	&	‘banana’	&	\tsc{Timor-Babar}, Timor		\\
\rcl	&	Luang	&	  {\hr}udi	&	‘banana’	&	\tsc{Timor-Babar}, SW Maluku		\\
\rcl	&	SE Babar	&	  {\hr}uty	&	‘banana’	&	\tsc{Timor-Babar}, SW Maluku		\\
\rcl	&	South Nuaulu	&	 {\hr}uri	&	‘banana’	&	\tsc{Ambon-Seram}		\\
\rcl	&	Geser	&	{\hr}fudi	&	‘banana’	&	STB, E. Seram		\\
\rcl	&	Yamdena	&	{\hr}funri	&	‘banana’	&	STB, Tanimbar		\\
\rcl	&	Arguni	&	{\hr}fud	&	‘banana’	&	STB, Bomberai Peninsula		\\
\rcl	&	Sekar	&	{\hr}fudi	&	‘banana’	&	STB, Bomberai Peninsula		\\
\rcl	&	Onin	&	{\hr}fugi	&	‘banana’	&	STB, Bomberai Peninsula (irr. \it{g})		\\
\rcl	&	Irarutu	&	{\hr}ɸud	&	‘banana’	&	STB, Bomberai Peninsula		\\
	&	Manam 	&	  {\hr}udi	&	‘banana’	&	\tsc{Oceanic}, PNG		\\
	&	Takia	&	{\hr}fud	&	‘banana’	&	\tsc{Oceanic}, PNG		\\
\lspbottomrule
\end{xltabular}\addtocounter{table}{-1}}
%\mbox{}\\}

Many Papuan languages reflect a very different set of cognates
for ‘banana’ (data in \qf{02-42} and \qf{02-43} from
\citet{DenhamDonohue2009}; supplemented
by the sources given in appendix \ref{ch:BanTer}).
The presence of “cognate” forms appears in unrelated language
families in the same region.

{\wg\begin{xltabular}{\linewidth}{>{\hspace{-\tabcolsep}}l<{\hspace{-\tabcolsep}}>{\hspace{-\tabcolsep}\enspace}llll}
\lsptoprule
\tbx{02-42}		&	Language	&	{\#}muku	&	‘banana’	&	Region		\\ \midrule \hhline{~}
%\endfirsthead
	%&	Language	&	{\#}muku	&	‘banana’	&	region		\\ \midrule
%\endhead
	&	Tehit	&	oga	&	‘banana’	&	West Papuan; W. Bird’s Head		\\
	&	Mor	&	moga	&	‘banana’	&	Isolate; Bomberai, Bird’s Head		\\
	&	Mbaham	&	muŋgʷó	&	‘banana’	&	TNG; W. Bomberai; Bird’s Head		\\
	&	Oirata	&	mu	&	‘banana’	&	TAP; Kisar, Maluku		\\
	&	Bunak	&	mok	&	‘banana’	&	TAP; Timor		\\
	&	Kula	&	maka	&	‘banana’	&	TAP; Alor		\\
	&	Wersing	&	məlul	&	‘banana’	&	TAP; Alor		\\
	&	Kamang	&	mooi {\tl} muui	&	‘banana’	&	TAP; Alor		\\
	&	Klon	&	mgol\footnotemark 	&	‘banana’	&	TAP; Alor		\\
	&	Pura	&	mol	&	‘banana’	&	TAP; Pantar		\\
	&	Teiwa	&	muxui	&	‘banana’	&	TAP; Pantar		\\
\lspbottomrule
\end{xltabular}\addtocounter{table}{-1}}
\footnotetext{
		Klon /mgol/ is phonetically [mŏgol]. It has an epenthetic
		“fast vowel” to break up consonant clusters.
		Comparatively such epenthetic vowels often mirror full vowels
		in related languages.}
		
Yet many Austronesian languages in eastern Indonesia
and Timor-Leste also have their main generic word for ‘banana’
in a form that is very close to the Papuan forms.
There are additional examples in \citet{DenhamDonohue2009}.

%{\tblr{@{}lllll}
{\wg\begin{xltabular}{\linewidth}{>{\hspace{-\tabcolsep}}l<{\hspace{-\tabcolsep}}>{\hspace{-\tabcolsep}\enspace}llll}
\lsptoprule
\tbx{02-43}		&	Language	&	{\#}muku	&	‘banana’	&	Region		\\ \midrule	\hhline{~}
\endfirsthead
\lsptoprule
\qf{02-43}	{\cwh}	&	Language	{\cwh}&	{\#}muku	{\cwh}&	‘banana’	{\cwh}&	region		{\cwh}\\ \midrule
\endhead
	&	Manggarai	&	muku	&	‘banana’	&	\tsc{Bima-Lembata}; west Flores\\ \hhline{~}
	&		&		&		&	(Pacar dialect has \it{punti})		\\
	&	Ngadha	&	muku	&	‘banana’	&	\tsc{Bima-Lembata}; central Flores		\\
	&	Sika	&	muʔu	&	‘banana’	&	\tsc{Bima-Lembata}; eastern Flores		\\
	&	Lewotobi	&	muko	&	‘banana’	&	\tsc{Bima-Lembata}; eastern Flores		\\
	&	C. Lembata	&	muku	&	‘banana’	&	\tsc{Bima-Lembata}; Lembata		\\
	&	Kedang	&	muʔu	&	‘banana’	&	\tsc{Bima-Lembata}; Lembata		\\
	&	Kambera	&	muku	&	‘banana sp.’	&	\tsc{Bima-Lembata}; east Sumba		\\
	&	Dhao	&	muʔu	&	‘banana’	&	\tsc{Bima-Lembata}; Havu-Dhao		\\
	&	Kemak	&	muu	&	‘banana’	&	\tsc{Central Timor}		\\
	&	S. Mambae	&	muu	&	‘banana’	&	\tsc{Central Timor}		\\
	&	Habun	&	muku	&	‘banana’	&	\tsc{Timor-Babar}; East Timor		\\
	&	Idate	&	muʔu	&	‘banana’	&	\tsc{Timor-Babar}; East Timor		\\
	&	Galolen	&	muʔu	&	‘banana’	&	\tsc{Timor-Babar}; East Timor		\\
	&	Dadu{\Q}a	&	muu	&	‘banana’	&	\tsc{Timor-Babar}; East Timor/Atauro		\\
	&	Raklungu	&	muʔu	&	‘banana’	&	\tsc{Timor-Babar}; Atauro		\\
	&	Tugun	&	muu	&	‘banana’	&	\tsc{Timor-Babar}; Wetar		\\
	&	Aputai	&	muu	&	‘banana’	&	\tsc{Timor-Babar}; Wetar		\\
	&	Kisar	&	muʔu	&	‘banana’	&	\tsc{Timor-Babar}; SW Maluku		\\
	&	Dobel	&	muʔu	&	‘banana’	&	\tsc{Aru}		\\
	&	Karey	&	mogo	&	‘banana’	&	\tsc{Aru}		\\
	&	Fordata	&	muʔu	&	‘banana’	&	STB; Tanimbar		\\
	&	Kei	&	muk	&	‘banana’	&	STB; Kei Islands		\\
	&	Kur	&	muk	&	‘banana’	&	STB; Teor-Kur		\\
	&	Bedoanas	&	maʔo	&	‘banana’	&	STB; Bomberai Peninsula		\\
	&	Erokwanas	&	maʔo	&	‘banana’	&	STB; Bomberai Peninsula		\\
	&	Uruangnirin	&	maŋgo	&	‘banana’	&	STB; Bomberai Peninsula		\\
\lspbottomrule
\end{xltabular}\addtocounter{table}{-1}}
%\mbox{}\\}

Note that the variation of forms in Papuan languages in \qf{02-42}
is much greater than the variation found in Austronesian languages in \qf{02-43}.
The relative stability of the forms in Austronesian languages
suggests that the Papuan term may have been borrowed only once
or twice very early in the Austronesian presence in Wallacea.
\citet[478,486]{UsherSchapper2022} more specifically link the words in Austronesian
languages associated with {\#}muku ‘banana’ ultimately
to Papuan Proto-Greater West Bomberai *maŋgʷVl ‘banana’
(where V is a vowel of uncertain quality),
and give this as the ultimate source for reflexes of **mugu
(or {\#}muku)
‘banana’ found in various Austronesian languages
throughout our target region.
So it appears that one variety of banana,
and also the word associated with it spread westward from the
New Guinea region.
This Papuan loan in Austronesian languages is clear evidence of contact.

\begin{figure}
	\centering
	\includegraphics[width=0.33\columnwidth]{02-19.jpg}
	\fcp{Banana (\it{Musa} sp.)}{02-19}
\end{figure}

As mentioned in \srf{02-01},
before the arrival of the Austronesian-speaking peoples,
the banana species \it{Musa acuminata ssp.
banksii} was already found on the New Guinea mainland
and was transported \it{westward} to the Lesser Sunda Islands
reaching as far as western Flores in what is modern day central Indonesia.
This is known from DNA studies
and is also reflected in the Papuan pseudo-etymon {\#}muku ‘banana’
and its cognates \citep{DenhamDonohue2009,DonohueDenham2009,DonohueDenham2010}.
The Austronesian-speaking peoples came into the region bringing
with them *punti ‘banana’
and its cognates, which originally referred to \it{Musa acuminata ssp. erans}.
The two subspecies of banana -- and their terms -- met in the middle -- what
we are calling the contact zone of
eastern Indonesia and Timor-Leste -- linguistic Wallacea.
Some Austronesian languages eventually adopted the Papuan word for ‘banana’.
In most of the modern day languages the terms above are their
generic term for ‘banana’,
except for languages of Sumba,
like Kambera, which denotes a certain type of banana.
Neighbouring related languages can have the counterpart term.
So Tetun has \it{hudi} ‘banana’ while nearby Galolen has \it{muʔu}
‘banana’ -- both Austronesian languages of the \tsc{Timor-Babar}
subgroup spoken on the Timor mainland.
And Kisar \it{muʔu} ‘banana’
and nearby Luang \it{udi} ‘banana’ are in the same \tsc{Southwest
Maluku} node of \tsc{Timor-Babar}.

\begin{figure}
	\includegraphics[width=\columnwidth]{muku.pdf}
	\caption
	[Distribution of reflexes of {\#}muku ‘banana’]
	{Distribution of reflexes of {\#}muku ‘banana’\footnotemark}
	\label{fig:02-20}
\end{figure}
\footnotetext{
Sources of data for \frf{02-20} are given in appendix \ref{ch:BanTer}.}

\citet[25]{DenhamDonohue2009} note: 

\begin{quote}
	“In the case of the distribution of *muku terms we know that
	the languages involved are not all related,
	and so a common ancestor cannot explain the distribution.
	The contemporary social situation in eastern Indonesia,
	and the fact that we are dealing with more than 50 widely-dispersed
	languages means we can discount the recent borrowing
	and the chance similarity explanations.
	Given that *muku does not occur in Austronesian languages away
	from the area of Papuan contact,
	a Papuan source must be posited.
	Despite the great time depth we are proposing,
	and the presence of the pseudo-cognate form in at least two unrelated
	language families, we believe that this lexical trace has persisted.
	[{\ldots}] Significantly, the geographical distribution of the
	term *muku, which is applied to cultivars today,
	extends from western New Guinea to western Flores.
	The term crosses a botanical boundary; it extends from the natural
	ranges of Musa and Callimusa (those formerly termed Australimusa) section species
	(see Figure 1) to regions where cultivars were introduced,
	ie, Timor and Flores.
	Taking into consideration the substratal effects of Papuan languages
	on modern eastern Austronesian languages,
	bananas referred to as *muku were spread westward over the current
	area in which the term is found by a pre-Austronesian culture
	with a maritime component.” \citep[25]{DenhamDonohue2009}
\end{quote}

The distribution of reflexes of {\#}muku ‘banana’ (see \frf{02-20})
has significant \it{methodological implications}.
They are found in unrelated Papuan language families.
They are also found widely in Austronesian languages that cannot
be assumed to have been in contact with each other since settling the area.
They are found in both Papuan
and Austronesian languages far to the west of the area of the original cultivars.
They are found along the southern chain of islands,
but not the northern chain of islands in our target region.
The \tsc{Seram-Tanimbar-Bomberai} subgroup has some nodes with
languages that have {\#}muku,
and other nodes that have none.
How did the plants get all the way to western Flores? How did
the words with such consistency in form
and meaning get into all these scattered languages
and subgroups, some of whose near neighbours maintained reflexes
of PMP *punti ‘banana’?
%We return to this issue in \crf{ch:MorLes}.
Fuller data for banana terms from 389 Austronesian
and Papuan languages are found in appendix \ref{ch:BanTer}.